% Created 2023-02-08 Wed 16:09
% Intended LaTeX compiler: pdflatex
\documentclass[11pt]{article}
\usepackage[utf8]{inputenc}
\usepackage[T1]{fontenc}
\usepackage{graphicx}
\usepackage{longtable}
\usepackage{wrapfig}
\usepackage{rotating}
\usepackage[normalem]{ulem}
\usepackage{amsmath}
\usepackage{amssymb}
\usepackage{capt-of}
\usepackage{hyperref}
\date{\today}
\title{greedy algorithm}
\hypersetup{
 pdfauthor={},
 pdftitle={greedy algorithm},
 pdfkeywords={},
 pdfsubject={},
 pdfcreator={Emacs 28.2 (Org mode 9.5.5)}, 
 pdflang={English}}
\begin{document}

\maketitle
\tableofcontents

the method of greediness proposes that in each step in an algorithm, where we are presented with multiple choices, the choice with the biggest instantaneous gain is picked\\
using this method we can write algorithm that, while relatively simple, arent usually optimal\\
\begin{question}
\uline{input}: an array \texttt{A} of positive numbers\\
\uline{output}: a subarray (non-consecutive) whose sum is maximal which doesnt contain two consecutive elements\\
\begin{answer}
the algorithm \(\textsc{Greedy1}(A)\) picks the biggest number in the input array, adds it to the total sum, and discards the two neighboring elements\\
\includesvg{/home/mahmooz/brain/out/1fIA2yy}\\

but is this algorithm correct?\\
\end{answer}
\end{question}
\end{document}